% Source: Wald GR, physical PDF p. 98 (printed p. 85).
% Coverage: Figure 4.2, spacelike hyperplane used to define radiated energy.
% Figures/tables/footnotes: Figure only; no tables or footnotes.
% Status: source-reviewed and compile-clean.
% Uncertainties: none.
\begingroup
\renewcommand{\thefigure}{4.2}
\begin{figure}[htbp]
  \centering
  \begin{tikzpicture}[scale=0.78,>=Stealth]
    \coordinate (A) at (-3.2,-1.25);
    \coordinate (B) at (2.8,-1.25);
    \coordinate (C) at (4.0,0.55);
    \coordinate (D) at (-2.0,0.55);
    \draw[thick] (A)--(B)--(C)--(D)--cycle;
    \node at (3.80,1.02) {\(\Sigma\)};
    \draw[->] (3.66,.90) to[bend right=14] (3.12,.48);

    \foreach \x/\y/\dx/\dy in {
      -2.35/2.55/1.15/-2.25,
      -.20/3.05/.25/-2.65,
      2.55/2.55/-1.15/-2.20,
      -1.25/-.55/-.70/-2.20,
      .65/-.65/.18/-2.35,
      2.25/-.38/.82/-2.22} {
      \draw[thick,decorate,decoration={snake,amplitude=2.8pt,segment length=10pt}]
        (\x,\y) -- ++(\dx,\dy);
    }
  \end{tikzpicture}
  \caption{A spacelike hyperplane, \(\Sigma\). The total energy contained in
  gravitational radiation is obtained by integrating \(t_{00}\) over \(\Sigma\).}
  \label{fig:4.2}
\end{figure}
\endgroup
